\chapter{Requirements Engineering}\label{ch:ch03}

Requirements engineering is the discipline of discovering, analysing, documenting,
and validating the services a system must provide and the constraints under which it
must operate~\cite{sommerville-se}. For APEX the exercise was neither a formal
customer-elicitation programme nor a speculative wish-list: the system was built to
serve a single, well-understood workflow --- running mixed-format examinations in a
university classroom --- and the requirements below were distilled from that workflow,
from the real feature set that the implementation ultimately shipped, and from the
honest constraints of a five-person graduation project. This chapter records those
requirements as they were actually realised, so that every functional requirement can
be traced to running code and every non-functional requirement to a measurable
property of the deployed system. Pressman's distinction between what a product must
\emph{do} (functional) and the qualities it must \emph{exhibit} (non-functional)
frames the presentation~\cite{pressman-se}.

The chapter is deliberately grounded rather than aspirational. Where a capability was
scoped out, or shipped with a known limitation, it is stated plainly; padding the
requirement set with features the code does not implement would defeat the purpose of a
traceability record.

\section{Stakeholders and Actors}
\label{sec:req-stakeholders}

A \emph{stakeholder} is any party with a legitimate interest in the system; an
\emph{actor} is a role that interacts directly with it across the system
boundary~\cite{sommerville-se}. APEX has few of the former and only two of the latter,
and that scarcity is itself a design decision worth stating.

\subsection{Stakeholders}

\begin{itemize}[leftmargin=1.4em]
  \item \textbf{Teachers} own the pedagogical content --- classes, exams, question
        banks --- and are the primary paying beneficiaries of automation: the system
        exists to remove the manual labour of grading code.
  \item \textbf{Students} are the subjects of assessment. Their overriding interest is
        that the timed sitting be fair, resilient to network and device failures, and
        transparent in how a score was reached.
  \item \textbf{The instructor / institution} (the project's academic supervisor and
        the Faculty of Engineering at Menoufia University) is a governance stakeholder:
        the system must be defensible as an engineering artefact and self-hostable so
        that no student data leaves institutional control.
  \item \textbf{The operator} --- whoever deploys and runs the instance --- has a stake
        in low operational cost, reproducible builds, and safe schema evolution. In a
        classroom deployment the operator and the teacher may well be the same person,
        which reinforces the requirement that operation be simple.
\end{itemize}

\subsection{Actors: two roles, no administrator}
\label{sec:req-actors}

APEX recognises exactly two interacting roles, encoded in the \texttt{role} column of
the \texttt{User} table as the string literals \texttt{teacher} and \texttt{student}.
\Cref{tab:actors} summarises them. There is deliberately \textbf{no administrator
role}. This is a genuine architectural boundary, not an omission: the authorization
middleware, \texttt{middleware.RequireRole(roles...)}, is generic and accepts any set
of role strings, but the shipped build only ever instantiates it with
\texttt{"teacher"} and \texttt{"student"}. The consequence is that every privileged
action must be justified by \emph{ownership of a specific resource} (a teacher may edit
only exams whose \texttt{teacher\_id} matches their own identity) rather than by a
blanket super-user grant. Removing the administrator removes an entire class of
privilege-escalation target and keeps the authorization story to two coarse roles plus
a per-row ownership check.

\begin{table}[htbp]
\centering
\caption{Primary actors of the APEX system.}
\label{tab:actors}
\small
\begin{tabularx}{\textwidth}{@{}L{2.4cm} L{3.2cm} X@{}}
\toprule
\textbf{Actor} & \textbf{Role literal} & \textbf{Goals and system interactions} \\
\midrule
Teacher & \texttt{teacher} &
Create and manage classes; author coding, MCQ, and written questions; maintain a
reusable question bank in folders; build, publish, and assign exams to classes;
auto- and manually grade submissions; post announcements; monitor results and class
statistics. \\
\addlinespace
Student & \texttt{student} &
Join classes by invite code; take timed exams in a browser code editor with autosave
and resume; run code against sample tests before submitting; review per-test results
and teacher feedback; receive notifications and reminders. \\
\bottomrule
\end{tabularx}
\end{table}

A handful of \emph{external} systems participate in the requirements without being
human actors: Judge0 (the code-execution sandbox reached over HTTP), an SMTP relay for
password-reset and reminder mail, and Google Identity Services for federated sign-in.
Each is optional-or-degradable --- absent an SMTP host, mail is logged to standard
output; absent a Google client ID, the federated endpoint returns HTTP~503 --- which is
itself a portability requirement captured in \Cref{sec:req-nfr}.

\section{User Stories}
\label{sec:req-stories}

Following the agile convention, requirements were first captured as user stories of the
form \emph{``As a \textless role\textgreater, I want \textless capability\textgreater\
so that \textless benefit\textgreater''}~\cite{pressman-se}. The stories below are the
ones that survived into shipped functionality; each maps to one or more functional
requirements in \Cref{sec:req-fr} and, through them, to concrete endpoints and handlers.

\subsection*{Teacher stories}

\begin{enumerate}[label=\textbf{ST-T\arabic*},leftmargin=3.2em]
  \item As a teacher, I want to create a class and share an invite code, so that
        students can enrol themselves without my manually maintaining a roster.
  \item As a teacher, I want to author coding, multiple-choice, and written questions
        with points, difficulty, hints, and time/memory limits, so that one exam can
        assess several kinds of skill in a single sitting.
  \item As a teacher, I want to keep reusable questions in a folder-organised bank, so
        that I can compose new exams from proven material instead of re-authoring it.
  \item As a teacher, I want exams to start as drafts and only become visible when I
        explicitly publish them (which requires setting a start time), so that a
        half-finished exam never leaks to students.
  \item As a teacher, I want to assign one exam to several classes at once, so that
        parallel sections sit the same paper.
  \item As a teacher, I want code answers graded automatically against my test cases,
        and written answers queued for me to grade with feedback and a score override,
        so that I spend my time only on the answers that genuinely need human judgement.
  \item As a teacher, I want to post class announcements that notify enrolled students,
        so that I can communicate without a separate channel.
  \item As a teacher, I want to preview an exam exactly as a student will see it and to
        withhold results until grading is complete, so that I control what is released
        and when.
\end{enumerate}

\subsection*{Student stories}

\begin{enumerate}[label=\textbf{ST-S\arabic*},leftmargin=3.2em]
  \item As a student, I want to join a class by typing an 8-character code, so that
        enrolment is instant and needs no teacher intervention.
  \item As a student, I want my in-progress answers saved on the server as I work, so
        that a dropped connection, a dead battery, or a switched device does not cost me
        my work.
  \item As a student, I want to resume the same attempt after an interruption rather
        than starting over, so that a crash mid-exam is recoverable.
  \item As a student, I want to run my code against sample tests and see the output
        before I submit, so that I get fast feedback while I still have time to fix it.
  \item As a student, I want to review, after grading, exactly which tests passed and
        any feedback my teacher left, so that the mark is transparent and I can learn
        from it.
\end{enumerate}

\subsection*{Platform stories}

\begin{enumerate}[label=\textbf{ST-P\arabic*},leftmargin=3.2em]
  \item As an operator, I want to stand up the whole system with one binary, a
        PostgreSQL database, and an external Judge0, so that a classroom instance is
        cheap and self-hostable.
  \item As an operator, I want the service to refuse to start with a weak signing
        secret and to evolve its schema safely, so that a misconfiguration cannot become
        a silent data-loss incident.
  \item As a user of either role, I want to be reminded of upcoming exams and notified
        when my work is graded, without duplicate messages if the server restarts.
\end{enumerate}

\section{Use-Case Model}
\label{sec:req-usecase}

The user stories condense into the use-case model of \cref{fig:usecase}, which shows the
principal interactions of each actor across the system boundary. The Teacher's use cases
cluster around \emph{authoring and administration} --- managing classes, building exams,
curating the question bank, assigning, grading, announcing, and inspecting results ---
while the Student's cluster around \emph{participation} --- joining, sitting the timed
exam, running and submitting code, reviewing results, and receiving notifications. The
asymmetry is expected: APEX automates the teacher's labour and constrains the student's
environment.

\begin{figure}[htbp]\centering
\resizebox{\textwidth}{!}{%
\begin{tikzpicture}[
    every node/.style={font=\small},
    actor/.style={
      inner sep=0pt, outer sep=2pt, minimum height=1.1cm, minimum width=0.7cm,
      align=center, text=apexblue, font=\small\bfseries
    },
    usecase/.style={
      draw=apexblue!70, thick, fill=apexlight, text=black!85,
      ellipse, align=center, inner sep=1.5pt, minimum width=3.2cm, minimum height=0.95cm,
      font=\scriptsize
    },
    link/.style={apexgrey, thin},
]

% ---- Stick-figure actor macro ----
\newcommand{\stickfig}[3]{% #1=name/anchor, #2=x, #3=y
  \coordinate (#1) at (#2,#3);
  \draw[apexblue, very thick] (#2,#3+0.55) circle (0.16);            % head
  \draw[apexblue, very thick] (#2,#3+0.39) -- (#2,#3-0.18);          % torso
  \draw[apexblue, very thick] (#2-0.30,#3+0.20) -- (#2+0.30,#3+0.20);% arms
  \draw[apexblue, very thick] (#2,#3-0.18) -- (#2-0.24,#3-0.60);     % left leg
  \draw[apexblue, very thick] (#2,#3-0.18) -- (#2+0.24,#3-0.60);     % right leg
}

% ================= System boundary =================
\node[
  draw=apexblue, very thick, rounded corners=32pt, fill=apexcyan!5,
  minimum width=9.2cm, minimum height=12.6cm,
] (boundary) at (0,0) {};
\node[anchor=north, apexblue, font=\bfseries] at (boundary.north) {\ APEX Classroom Portal\ };

% ================= Teacher use cases (left column) =================
\node[usecase] (uc-classes)  at (-2.0, 5.1) {Manage Classes};
\node[usecase] (uc-build)    at (-2.0, 3.6) {Build Exams};
\node[usecase] (uc-bank)     at (-2.0, 2.1) {Manage\\Question Bank};
\node[usecase] (uc-assign)   at (-2.0, 0.6) {Assign Exams};
\node[usecase] (uc-grade)    at (-2.0,-0.9) {Grade\\Written Answers};
\node[usecase] (uc-announce) at (-2.0,-2.4) {Post Announcements};
\node[usecase] (uc-stats)    at (-2.0,-3.9) {View Results / Stats};

% ================= Student use cases (right column) =================
\node[usecase] (uc-join)     at ( 2.0, 4.35) {Join Class};
\node[usecase] (uc-take)     at ( 2.0, 2.85) {Take Timed Exam};
\node[usecase] (uc-run)      at ( 2.0, 1.35) {Run / Submit Code};
\node[usecase] (uc-review)   at ( 2.0,-0.15) {Review Results};
\node[usecase] (uc-notify)   at ( 2.0,-1.65) {Receive\\Notifications};

% ================= Actors =================
\stickfig{teacher}{-6.6}{0.6}
\node[actor] at (-6.6,-0.5) {Teacher};

\stickfig{student}{6.6}{0.6}
\node[actor] at (6.6,-0.5) {Student};

% ================= Associations: Teacher =================
\draw[link] (teacher) -- (uc-classes.west);
\draw[link] (teacher) -- (uc-build.west);
\draw[link] (teacher) -- (uc-bank.west);
\draw[link] (teacher) -- (uc-assign.west);
\draw[link] (teacher) -- (uc-grade.west);
\draw[link] (teacher) -- (uc-announce.west);
\draw[link] (teacher) -- (uc-stats.west);

% ================= Associations: Student =================
\draw[link] (student) -- (uc-join.east);
\draw[link] (student) -- (uc-take.east);
\draw[link] (student) -- (uc-run.east);
\draw[link] (student) -- (uc-review.east);
\draw[link] (student) -- (uc-notify.east);

\end{tikzpicture}%
}
\caption{Principal use cases for the Teacher and Student actors.}\label{fig:usecase}
\end{figure}


Two properties of \cref{fig:usecase} deserve comment. First, no use case crosses from
one actor to the other; the ownership model in \Cref{sec:req-actors} enforces that a
teacher never ``takes'' an exam and a student never ``grades'' one. Second, several
teacher use cases (\textsl{Grade Written Answers}, \textsl{View Results / Stats}) and
the student's \textsl{Review Results} are two ends of the same underlying data ---
the attempt aggregate --- so although they are separate use cases they are bound by a
shared invariant: a student sees a result only after, and exactly as, the teacher's
grading produced it.

\section{Functional Requirements}
\label{sec:req-fr}

Functional requirements state what the system must do. They are grouped below by the
party they primarily serve --- Teacher, Student, and Platform --- and each is given a
stable identifier (\texttt{FR-T\emph{n}}, \texttt{FR-S\emph{n}}, \texttt{FR-P\emph{n}})
so that \Cref{sec:req-trace} can trace it to an implementing component. Every
requirement listed corresponds to functionality that is present in the shipped system;
the wording is drawn from the real capabilities enumerated in the project's feature
inventory.

\subsection{Teacher requirements}
\label{sec:req-fr-teacher}

\begin{table}[htbp]
\centering
\caption{Functional requirements serving the Teacher actor.}
\label{tab:fr-teacher}
\small
\begin{tabularx}{\textwidth}{@{}L{1.6cm} X@{}}
\toprule
\textbf{ID} & \textbf{Requirement} \\
\midrule
FR-T1 & The system shall let a teacher create, edit, and delete classes, each with a
name, section, optional cover image, and an automatically generated unique 8-character
invite code, and shall list the class roster and permit removal of members. \\
FR-T2 & The system shall provide an exam builder for authoring questions of type
\texttt{coding}, \texttt{mcq}, and \texttt{written}, each with points, difficulty,
optional hints, image attachments, and tags; coding questions shall additionally carry
a per-question time limit (default \texttt{2000}\,ms) and memory limit (default
\texttt{262144}\,KB) and a set of test cases. \\
FR-T3 & The system shall maintain a per-teacher question bank of reusable questions
(\texttt{is\_bank = true}), organisable into folders, so that questions can be composed
into exams independently of any single exam. \\
FR-T4 & The system shall distinguish draft from published exams and shall refuse to
publish (clear \texttt{is\_draft}) an exam whose \texttt{start\_time} is unset, so that
only complete, scheduled exams become visible to students. \\
FR-T5 & The system shall let a teacher assign one exam to one or many classes, rejecting
any assignment that references a class the teacher does not own. \\
FR-T6 & The system shall configure per-exam options: shuffle question order, show
results after submission, and a passing score. \\
FR-T7 & The system shall automatically grade coding submissions against their test cases
and MCQ submissions by set-equality of the selected against the correct options, and
shall queue written submissions for manual grading. \\
FR-T8 & The system shall provide a pending-grading queue from which a teacher grades
written (and any unconfigured MCQ) answers, supplying a score, a status, and optional
free-text feedback, after which the affected attempt's aggregate score is recomputed. \\
FR-T9 & The system shall let a teacher post per-class announcements with attachments,
fanning out an in-app notification to every enrolled student. \\
FR-T10 & The system shall present a teacher dashboard, a results explorer with per-exam
outcomes, and per-class statistics. \\
FR-T11 & The system shall let a teacher preview an exam exactly as a student will see it,
and shall optionally block grade announcement for a class until all manual grading for it
is complete. \\
FR-T12 & The system shall, when a teacher destructively edits a published exam, wipe the
affected attempts, submissions, and test results in a single transaction, stamp
\texttt{reset\_at}, un-announce the class grades, and notify affected students. \\
\bottomrule
\end{tabularx}
\end{table}

\subsection{Student requirements}
\label{sec:req-fr-student}

\begin{table}[htbp]
\centering
\caption{Functional requirements serving the Student actor.}
\label{tab:fr-student}
\small
\begin{tabularx}{\textwidth}{@{}L{1.6cm} X@{}}
\toprule
\textbf{ID} & \textbf{Requirement} \\
\midrule
FR-S1 & The system shall let a student join a class by entering its 8-character invite
code, leave a class, and list the classes and assigned exams available to them. \\
FR-S2 & The system shall let a student start a timed attempt at a published, in-window
exam, creating at most one attempt per \texttt{(student, exam)} pair; a repeated start
shall resume the existing attempt rather than create a second one. \\
FR-S3 & The system shall autosave a student's in-progress answers to the server (\texttt{PUT
/student/exams/:id/autosave}) and rehydrate them on resume, so that a crash or device
change does not lose work. \\
FR-S4 & The system shall optionally present an exam's questions in shuffled order,
consistently within a single attempt. \\
FR-S5 & The system shall provide a browser code editor supporting six languages (Python,
JavaScript, TypeScript, Java, C, C++) with a run-against-sample-tests action that returns
output, time, and memory without persisting a graded submission. \\
FR-S6 & The system shall let a student submit an exam, after which no further submission
for that attempt is accepted, and shall grade the submission server-side. \\
FR-S7 & The system shall let a student review each attempt's result --- the aggregate
score, the per-test breakdown for coding questions, and any teacher feedback --- subject
to the teacher's results-release settings. \\
FR-S8 & The system shall present the student's notifications, unread count, and a help
page, and shall let the student manage a profile and notification preferences. \\
\bottomrule
\end{tabularx}
\end{table}

\subsection{Platform requirements}
\label{sec:req-fr-platform}

\begin{table}[htbp]
\centering
\caption{Platform-wide functional requirements (cross-cutting both roles).}
\label{tab:fr-platform}
\small
\begin{tabularx}{\textwidth}{@{}L{1.6cm} X@{}}
\toprule
\textbf{ID} & \textbf{Requirement} \\
\midrule
FR-P1 & The system shall authenticate users by email and password (bcrypt-hashed) and
issue a signed session token valid for 24 hours, and shall support federated sign-in via
Google Identity Services when configured. \\
FR-P2 & The system shall support password reset by email using a single-use, SHA-256-hashed
token with a one-hour expiry, and shall not reveal whether an email address is registered. \\
FR-P3 & The system shall let a user permanently delete their own account, cascading the
deletion across their dependent records in a single transaction. \\
FR-P4 & The system shall send in-app and email reminders for upcoming exams --- in-app at
roughly one hour before, email at one hour and again at start --- deduplicated by
per-exam timestamp columns so that a process restart cannot re-send a reminder. \\
FR-P5 & The system shall notify a student exactly once when their exam has been graded,
deduplicated by a per-attempt flag so that repeated finalisation does not repeat the
notification. \\
FR-P6 & The system shall expose an unauthenticated liveness endpoint (\texttt{GET
/healthz}) and a self-probe (\texttt{--healthcheck}) so that a shell-less container image
can be health-checked. \\
FR-P7 & The system shall evolve its database schema through idempotent, ordered
migrations, and in production shall make a versioned migration track the sole schema
authority (with GORM auto-migration disabled). \\
\bottomrule
\end{tabularx}
\end{table}

A note on grading semantics belongs with the functional requirements because it governs
observable behaviour. The aggregate score of an attempt (FR-S7, FR-T8) is not an
incrementally mutated counter but a value recomputed from scratch over the exam's
\emph{full} problem list, so a skipped question counts zero in the denominator, while a
\texttt{pending\_review} written answer is excluded from both numerator and denominator
until a human grades it, so an in-flight manual grade never transiently depresses the
visible total. This behaviour is the functional contract the traceability table binds to
the \texttt{maybeFinalizeAttempt} aggregator.

\section{Non-Functional Requirements}
\label{sec:req-nfr}

Non-functional requirements constrain \emph{how well} the system performs its functions.
To be useful they must be measurable~\cite{sommerville-se}; each requirement below is
therefore stated with a concrete, checkable criterion drawn from the actual
configuration and design of the system, and each is given an identifier
\texttt{NFR-\emph{n}}.

\subsection{Security}
\label{sec:req-nfr-security}

\begin{itemize}[leftmargin=1.4em]
  \item \textbf{NFR-1 (authentication strength).} Passwords shall be stored only as
        bcrypt hashes at \texttt{bcrypt.DefaultCost} (cost 10) and never returned in any
        API response; session tokens shall be HS256-signed JWTs with a 24-hour expiry.
  \item \textbf{NFR-2 (fail-fast secret).} The service shall refuse to start (via
        \texttt{log.Fatal}) if \texttt{JWT\_SECRET} is empty, shorter than 32 characters,
        or the placeholder literal \texttt{dev-secret-change-in-prod}. This is a boolean,
        boot-time-checkable criterion.
  \item \textbf{NFR-3 (defence in depth).} Every protected request shall pass three
        independent authorization gates that each fail closed: JWT verification (401 on
        any doubt), role check (403 on mismatch), and per-row ownership check in the
        handler; a failure in any one gate shall not open the others.
  \item \textbf{NFR-4 (input and injection safety).} The system shall parameterise all
        database access through the ORM (mitigating \texttt{CWE-89}~\cite{cwe-89}),
        restrict cross-origin access to a configured allow-list, and never execute
        student code in the application process --- all execution shall be delegated to
        the Judge0 \texttt{isolate} sandbox~\cite{judge0,isolate}.
  \item \textbf{NFR-5 (honest limits).} The requirement set explicitly records the known
        gaps as accepted risk for the thesis snapshot: authentication endpoints are not
        rate-limited, the password minimum is six characters, and the session token lives
        in browser \texttt{localStorage}. Each is a documented decision, to be hardened
        before public exposure, not a hidden defect.
\end{itemize}

The security requirements are aligned with the OWASP Top 10 categories for broken access
control and identification failures~\cite{owasp-top10}; \Cref{ch:ch03} states them,
\cref{fig:threat} and the security chapter analyse them.

\subsection{Performance and scalability}
\label{sec:req-nfr-perf}

\begin{itemize}[leftmargin=1.4em]
  \item \textbf{NFR-6 (responsiveness of submit).} Exam submission shall return to the
        student without waiting for grading: grading is dispatched to background workers
        so the HTTP response carries only the attempt and submission identifiers. The
        measurable criterion is that the submit response latency is independent of the
        number or runtime of test cases.
  \item \textbf{NFR-7 (bounded execution).} Each code execution shall run under the
        per-problem time limit (default \texttt{2000}\,ms) and memory limit (default
        \texttt{262144}\,KB) enforced by the sandbox, and the grading HTTP client shall
        apply a 30-second ceiling per Judge0 call, so a pathological submission cannot
        occupy resources unbounded.
  \item \textbf{NFR-8 (connection pooling).} The database layer shall bound its
        connection pool --- 25 max-open, 10 max-idle connections, with 5-minute idle and
        30-minute maximum connection lifetimes --- so that concurrent load degrades
        gracefully rather than exhausting the database.
  \item \textbf{NFR-9 (indexed hot paths).} Correctness- and performance-critical queries
        shall be backed by indexes: unique composite indexes on
        \texttt{ClassMember(class\_id, user\_id)}, \texttt{ExamAttempt(user\_id, exam\_id)},
        and \texttt{ExamClass(exam\_id, class\_id)} enforce their one-row invariants, and
        \texttt{exams.reset\_at} is indexed to support the take-screen cache-bust query.
\end{itemize}

The target scale is a classroom, not a MOOC: tens of concurrent students per exam, not
tens of thousands. The pooling and single-binary choices are sized for that target and
are stated as such rather than claimed to be web-scale.

\subsection{Usability}
\label{sec:req-nfr-usability}

\begin{itemize}[leftmargin=1.4em]
  \item \textbf{NFR-10 (heuristic conformance).} The user interface shall conform to
        Nielsen's usability heuristics~\cite{nielsen-heuristics}, in particular
        \emph{visibility of system status} (a submission's state --- \texttt{pending},
        \texttt{running}, graded --- is always shown; live sample-test feedback reports
        pass/fail before submit) and \emph{error prevention} (a draft exam cannot be
        published without a start time; a submitted attempt cannot be resubmitted).
  \item \textbf{NFR-11 (recoverability).} The interface shall protect against user and
        environment error by never silently discarding work: in-progress answers are
        autosaved server-side and an interrupted attempt resumes where it left off,
        satisfying the ``user control and freedom'' heuristic.
  \item \textbf{NFR-12 (measured satisfaction).} Perceived usability shall be evaluable
        with the System Usability Scale~\cite{brooke1996sus}, a ten-item instrument
        yielding a 0--100 score; adopting a standard instrument makes the usability claim
        falsifiable rather than anecdotal. The evaluation is reported in the results
        chapter.
\end{itemize}

\subsection{Portability and self-hostability}
\label{sec:req-nfr-portability}

\begin{itemize}[leftmargin=1.4em]
  \item \textbf{NFR-13 (single-artefact deploy).} The backend shall compile to a single
        statically linked binary with no runtime dependencies beyond PostgreSQL and a
        reachable Judge0, and shall ship in a distroless container image running as a
        non-root user.
  \item \textbf{NFR-14 (twelve-factor configuration).} All configuration shall be read
        from environment variables with sensible defaults; no configuration shall be
        compiled in, so the same image runs in development and production unchanged.
  \item \textbf{NFR-15 (graceful degradation of externals).} Optional external services
        shall degrade rather than fail: with \texttt{SMTP\_HOST} unset, mail is logged to
        standard output; with \texttt{GOOGLE\_CLIENT\_ID} unset, the federated sign-in
        endpoint returns HTTP~503 while password sign-in continues to work; the default
        Judge0 URL points at a public instance so first-run needs no local sandbox.
  \item \textbf{NFR-16 (managed-database compatibility).} The system shall accept a full
        \texttt{DATABASE\_URL} (taking precedence over discrete \texttt{DB\_*} variables),
        making it deployable against managed PostgreSQL providers without code change.
\end{itemize}

\subsection{Maintainability and reliability}
\label{sec:req-nfr-maint}

\begin{itemize}[leftmargin=1.4em]
  \item \textbf{NFR-17 (layered structure).} The backend shall be organised into four
        request-processing layers (transport, handlers, domain services, persistence)
        with feature-scoped packages, so that a change is localised to one package.
  \item \textbf{NFR-18 (safe schema evolution).} Schema changes shall follow the
        add-nullable $\rightarrow$ backfill $\rightarrow$ constrain discipline and shall
        never rely on an ORM DDL default to backfill a populated table; production shall
        run a versioned migration track as the sole schema authority. This requirement
        was forced by a real incident (an ORM default silently hid every exam from
        students) and is now enforced by project rule.
  \item \textbf{NFR-19 (continuous integration).} Every change shall pass an automated
        pipeline --- backend vet, test, and lint; frontend lint, build, and test; and a
        Docker smoke build --- before merge, so that the mainline remains buildable.
  \item \textbf{NFR-20 (idempotent background work).} Reminder dispatch and grade
        notification shall be idempotent across process restarts, using database columns
        rather than in-memory state for deduplication, so that reliability does not depend
        on process uptime.
\end{itemize}

APEX's configuration and deployment posture follow the twelve-factor
guidance~\cite{twelvefactor} and the distroless-image practice~\cite{distroless}; the
maintainability requirements reflect that lineage.

\section{Requirements Traceability}
\label{sec:req-trace}

Traceability links each requirement to the artefact that satisfies it, so that a reader
--- or an examiner --- can verify the claim against the source rather than take it on
trust~\cite{sommerville-se}. \Cref{tab:trace} maps every functional requirement to the
principal component and endpoint that realise it. The non-functional requirements trace
to cross-cutting mechanisms already named in \Cref{sec:req-nfr}: NFR-1--5 to the
authentication and authorization middleware and the Judge0 sandbox, NFR-6--9 to the
grading dispatch and the database pool and indexes, NFR-10--12 to the SPA and its
autosave/resume flow, NFR-13--16 to the Docker and configuration layer, and NFR-17--20 to
the package structure, migration pipeline, CI workflow, and reminder scheduler.

\begin{table}[htbp]
\centering
\caption{Functional-requirement traceability to implementing components.}
\label{tab:trace}
\small
\begin{tabularx}{\textwidth}{@{}L{1.4cm} L{4.2cm} X@{}}
\toprule
\textbf{Req.} & \textbf{Component / package} & \textbf{Principal endpoint(s)} \\
\midrule
FR-T1  & \texttt{internal/class}        & \texttt{POST/GET/PUT/DELETE /classes}, \texttt{DELETE /classes/:id/members/:userId} \\
FR-T2  & \texttt{internal/problem}, \texttt{internal/testcase} & \texttt{POST /exams/:id/problems}, \texttt{POST /problems/:id/test-cases} \\
FR-T3  & \texttt{internal/problem}, \texttt{internal/folder}   & \texttt{POST /problems/bank}, \texttt{GET/POST/PUT/DELETE /folders} \\
FR-T4  & \texttt{internal/exam}         & \texttt{PUT /exams/:id} (publish requires \texttt{start\_time}) \\
FR-T5  & \texttt{internal/exam}         & \texttt{POST /exams/:id/assign} \\
FR-T6  & \texttt{internal/exam}         & \texttt{POST/PUT /exams} \\
FR-T7  & \texttt{internal/judge0}       & driven by \texttt{POST /student/exams/:id/submit} \\
FR-T8  & \texttt{internal/teacher}, \texttt{internal/submission} & \texttt{GET /teacher/grading/pending}, \texttt{PUT /submissions/:id/grade} \\
FR-T9  & \texttt{internal/announcement}, \texttt{internal/notification} & \texttt{POST /classes/:id/announcements} \\
FR-T10 & \texttt{internal/teacher}, \texttt{internal/class}    & \texttt{GET /teacher/dashboard}, \texttt{GET /exams/:id/results}, \texttt{GET /classes/:id/stats} \\
FR-T11 & \texttt{internal/exam}, \texttt{internal/class}       & exam preview (SPA route), \texttt{grades\_announced} gating \\
FR-T12 & \texttt{internal/exam} (\texttt{reset.go}) & destructive edit triggers \texttt{reset\_at} \\
\addlinespace
FR-S1  & \texttt{internal/student}      & \texttt{POST /student/classes/join}, \texttt{DELETE /student/classes/:id} \\
FR-S2  & \texttt{internal/exam}         & \texttt{POST /student/exams/:id/start} (idempotent) \\
FR-S3  & \texttt{internal/exam}         & \texttt{PUT /student/exams/:id/autosave} \\
FR-S4  & \texttt{internal/exam}         & \texttt{shuffle\_questions} on the exam \\
FR-S5  & \texttt{internal/execute}, \texttt{internal/submission} & \texttt{POST /execute}, \texttt{POST /submissions/run} \\
FR-S6  & \texttt{internal/exam}, \texttt{internal/judge0}      & \texttt{POST /student/exams/:id/submit} \\
FR-S7  & \texttt{internal/submission}, \texttt{internal/student} & \texttt{GET /submissions/:id}, \texttt{GET /attempts/mine} \\
FR-S8  & \texttt{internal/notification}, \texttt{internal/profile} & \texttt{GET /notifications}, \texttt{GET/PUT /profile} \\
\addlinespace
FR-P1  & \texttt{internal/auth}, \texttt{internal/middleware} & \texttt{POST /auth/login}, \texttt{/auth/signup}, \texttt{/auth/google} \\
FR-P2  & \texttt{internal/auth}, \texttt{internal/email}      & \texttt{POST /auth/forgot-password}, \texttt{/auth/reset-password} \\
FR-P3  & \texttt{internal/auth}         & \texttt{DELETE /auth/account} \\
FR-P4  & \texttt{internal/reminder}, \texttt{internal/email}  & background scheduler (1-min tick) \\
FR-P5  & \texttt{internal/judge0}       & \texttt{maybeFinalizeAttempt} + \texttt{graded\_notified} \\
FR-P6  & \texttt{main.go}               & \texttt{GET /healthz}, \texttt{--healthcheck} \\
FR-P7  & \texttt{internal/database}, \texttt{cmd/migrate}     & migration passes / versioned SQL track \\
\bottomrule
\end{tabularx}
\end{table}

\section{Constraints and Assumptions}
\label{sec:req-constraints}

Finally, the requirements are bounded by explicit constraints and assumptions that shape
what the system may assume of its environment:

\begin{itemize}[leftmargin=1.4em]
  \item \textbf{Technology constraints.} The backend is Go with the Gin framework over
        PostgreSQL~16; code execution is delegated to Judge0~CE. These are fixed choices,
        not requirements to be re-derived, and are justified in the design and background
        chapters.
  \item \textbf{Deployment assumption.} A reachable Judge0 instance exists; the default
        points at the rate-limited public \texttt{ce.judge0.com}, and a real classroom
        deployment is assumed to run a self-hosted instance.
  \item \textbf{Scale assumption.} The concurrency target is classroom-sized; the system
        does not attempt horizontal scale-out, distributed grading queues, or multi-region
        deployment.
  \item \textbf{Trust assumption.} Students take exams on their own machines; APEX does
        not perform proctoring or lockdown, and treats the browser as untrusted --- which
        is precisely why all grading and all answer state live on the server.
\end{itemize}

These constraints are not deficiencies to apologise for; stating them is what makes the
functional and non-functional requirements above honest and achievable within the scope
of a graduation project. The chapters that follow show how each requirement was designed
(\Cref{ch:ch03} onward), implemented, secured, and evaluated.
